\documentclass[conference]{IEEEtran}
\IEEEoverridecommandlockouts
\usepackage{cite}
\usepackage{amsmath,amssymb,amsfonts}
\usepackage{algorithmic}
\usepackage{graphicx}
\usepackage{textcomp}
\usepackage{xcolor}
\usepackage{type1cm}
\usepackage{lettrine}
\usepackage{subfig}
\usepackage{fancyhdr} % fancy package

%%%%%%%%%%%%%%%%%%%%%%%%%%%%%%%%%%%%%%%%%%%%%%%FooterLift
\makeatletter
\def\ps@headings{% default to standard twoside headers, no footers
% will change later if the mode requires otherwise
\def\@oddhead{\@IEEEheaderstyle\hfil}%
\def\@evenhead{\@IEEEheaderstyle\hfil\hbox{}}% not used, oneside
\def\@oddfoot{\@IEEEheaderstyle\hfil\thepage}%
\def\@evenfoot{\@IEEEheaderstyle\thepage\hfil\hbox{}}% not used, oneside
}
\def\ps@IEEEtitlepagestyle{% default title page headers, no footers
\def\@oddhead{\hfil2021 22nd International Middle East Power Systems Conference (MEPCON), Assiut University, Egypt\hfil}%
\let\@evenhead\@empty
\def\@oddfoot{\footnotesize 978-1-6654-1998-7/21/\$31.00~\copyright 2021 IEEE\hfill\thepage}%
\let\@evenfoot\@empty
}
\makeatother
%%%%%%%%%%%%%%%%%%%%%%%%%%%%%%%%%%%%%%%%%%%%%%%FooterRight

\pagestyle{headings}
\setcounter{page}{77}
\thispagestyle{IEEEtitlepagestyle}

%%%%%%%orcid%%%%%%%%%%%
\usepackage{scalerel}
\usepackage{tikz}
\usetikzlibrary{svg.path}

\definecolor{orcidlogocol}{HTML}{A6CE39}
\tikzset{
	orcidlogo/.pic={
		\fill[orcidlogocol] svg{M256,128c0,70.7-57.3,128-128,128C57.3,256,0,198.7,0,128C0,57.3,57.3,0,128,0C198.7,0,256,57.3,256,128z};
		\fill[white] svg{M86.3,186.2H70.9V79.1h15.4v48.4V186.2z}
		svg{M108.9,79.1h41.6c39.6,0,57,28.3,57,53.6c0,27.5-21.5,53.6-56.8,53.6h-41.8V79.1z M124.3,172.4h24.5c34.9,0,42.9-26.5,42.9-39.7c0-21.5-13.7-39.7-43.7-39.7h-23.7V172.4z}
		svg{M88.7,56.8c0,5.5-4.5,10.1-10.1,10.1c-5.6,0-10.1-4.6-10.1-10.1c0-5.6,4.5-10.1,10.1-10.1C84.2,46.7,88.7,51.3,88.7,56.8z};
	}
}

\newcommand\orcidicon[1]{\href{https://orcid.org/#1}{\mbox{\scalerel*{
				\begin{tikzpicture}[yscale=-1,transform shape]
					\pic{orcidlogo};
				\end{tikzpicture}
			}{|}}}}

\usepackage[hidelinks]{hyperref}
\DeclareRobustCommand*{\IEEEauthorrefmark}[1]{%
	\raisebox{0pt}[0pt][0pt]{\textsuperscript{\footnotesize #1}}%
}


\begin{document}

\title{Performance Enhancement of Wave Energy Converters Integrated with DC Grids Using Energy Storage Systems}
\author{\IEEEauthorblockN{Mohamed G. Hussien \orcidicon{0000-0001-6410-5678}}
\IEEEauthorblockA{\textit{Department of Electrical Power and Machines Engineering} \\
\textit{Faculty of Engineering, Tanta University}\\
 Tanta, Egypt  \\
mohamed.hussien3@f-eng.tanta.edu.eg}
\and
\IEEEauthorblockN{Ahmed M. Elkholy \orcidicon{0000-0002-1834-1175}}
\IEEEauthorblockA{\textit{\small Department of Electrical Power and Machines Engineering} \\
\textit{Faculty of Engineering, Tanta University}\\
 Tanta, Egypt  \\
ahmed\_elkholy@f-eng.tanta.edu.eg}
\and
\IEEEauthorblockN{Diaa-Eldin A. Mansour \orcidicon{0000-0002-3894-4299}}
\IEEEauthorblockA{\textit{\small Department of Electrical Power and Machines Engineering} \\
	\textit{Faculty of Engineering, Tanta University}\\
	Tanta, Egypt  \\
	mansour@f-eng.tanta.edu.eg}
\and
\IEEEauthorblockN{Mahmoud M. Akl \orcidicon{0000-0002-6558-2449}}
\IEEEauthorblockA{\textit{Department of Electrical Power and Machines Engineering} \\
	\textit{Faculty of Engineering, Tanta University}\\
	Tanta, Egypt  \\
	mahmoud\_akl@f-eng.tanta.edu.eg}
\and
\IEEEauthorblockN{Mohamed Numair \orcidicon{0000-0001-9271-6759}}
\IEEEauthorblockA{\textit{Department of Electrical Power and Machines Engineering} \\
	\textit{Faculty of Engineering, Tanta University}\\
	Tanta, Egypt  \\
	mnumair@asu.edu}
\and
\IEEEauthorblockN{Ismail A. Soliman \orcidicon{0000-0002-0397-3719}}
\IEEEauthorblockA{\textit{Department of Electrical Power and Machines Engineering} \\
	\textit{Faculty of Engineering, Tanta University}\\
	Tanta, Egypt  \\
	esmail\_soliman@f-eng.tanta.edu.eg}
}


\maketitle
\
\begin{abstract}
Wave energy has a great potential as a renewable energy source. However, the intermittent nature of wave profile causes fluctuations in the output power extracted from wave energy converters. This paper proposes to smooth the output power extracted from wave energy using battery energy storage system and to integrate this power with DC systems. First, oscillating water column used to extract wave energy is described. It is based on reference model 3, which is referred as wave power buoy. Then, the design considerations of wave power buoy with permanent magnet linear generator are discussed and the Egyptian western coastline is taken as a case study. The ocean wave simulation is performed using the WecSim toolbox developed in MATLAB/SIMULINK considering the real irregular wave action. The output power from the linear generator is demonstrated. After that, the wave energy power generation with battery energy storage system is described. The system performance is evaluated through various simulation results. Finally, supercapacitor is used with battery for further enhancement in the performance of the system.
\end{abstract}

\begin{IEEEkeywords}
Wave Energy, PM linear generator, Egypt costs, Dc micogrid, WecSim
\end{IEEEkeywords}

\section{Introduction}
\lettrine{T}he limited fossil fuel with the corresponding increase in its cost forced electrical utilities to move forward renewable energy sources \cite{elkholy2020general,elkholy2018advanced}. In this regard, wave energy represented one of the attractive renewable energy sources due to its cleanness and availability on a wide scale. The wave energy potential has been assessed by many parties \cite{1,2,eslkoly2017islanding}. In \cite{1}, the assessment was done in the Gulf of Mexico and Caribbean over a period of 30 years and the sites having maximum energy potential could be identified. In \cite{2}, current methods used for characterizing wave energy resources and identifying expected generation were revised. Table.~\ref{tabintro} presents the wave energy potential for different locations all over the world \cite{3} Reference \cite{3}  presents the wave energy potential for different locations all over the world. In order to extract wave energy, usually wave energy converters through linear generators are used due to its effectiveness, reduced complexity and  less maintenance time due to the less moving parts number comparing to other technologies. There are various technologies of wave energy converters such as Oscillating Water Column, Archimedes wave swing, Oscillating Bodies, and so on \cite{3}.


\begin{table}[htbp]
	\caption{Wave energy potential for different locations all over the world \cite{3}}
	\begin{center}
		\begin{tabular}{|c|c|}
		 \hline
		 \hline
		 \\[-1em]
		 Region&	Wave energy potential \\
		 \hline
		 \hline
		 \\[-1em]
		 Mediterranean Sea and Atlantic Archipelagos&	$1300\ TWh/y$\\
		 \hline
		 \\[-1em]
		 Central America&	$1500\ TWh/y$ \\
		 \hline
		 \\[-1em]
		 Western and Northern Europe& 	$2800\ TWh/y$ \\
		 \hline
		 \\[-1em]
		 Africa&	$3500\ TWh/y$ \\
		 \hline
		 \\[-1em]
		 North America and Greenland& 	$4000\ TWh/y$ \\
		 \hline
		 \\[-1em]
		 South America&	$4600\ TWh/y$ \\
		 \hline
		 \\[-1em]
		 Australia, New Zealand and Pacific Islands& 	$5600\ TWh/y$\\
		 \hline
		 \\[-1em]
		 Asia& 	$6200\ TWh/y$\\
		 \hline
		 \hline
		\end{tabular}
		\label{tabintro}
	\end{center}
\end{table}

Grid integration of wave energy converters has gained a great interest in recent years \cite{4}. However, there are several problems that are encountered when integrating wave energy converters. The first problem is synchronization issues when integrating with AC grids. The second problem of integrating wave energy converters is the intermittent nature of wave profile causing fluctuations in the output power from wave energy converters. These power fluctuations create major operational concerns including voltage and frequency deviations as well as instabilities. As a result, this paper aims to mitigate these problems through using DC transmission and battery energy storage for grid integration of wave energy converters. DC transmission provides several advantages compared to AC one, such as simple control, flexible operation, absence of synchronization issues, and easiness for integrating energy storage devices \cite{5}. Accordingly, many researches have been done regarding integration of renewable energy sources with DC systems \cite{6,7}. On the other hand, battery energy storage will enable smoothing the output power from wave energy converters.

The paper is organized as follows. First, the wave power buoy is described along with its design with the linear generator. Next, the usage of adopted wave energy converter in the western Egyptian coastline is investigated as a case study. Finally, the integration of battery energy storage with wave energy converter is presented including modeling, control, and obtained results.

\section{RM3 Description}
The reference model 3 (RM3) is wave power buoy \cite{neary2014methodology}; this device is designed by "ocean power technology's power buoy". It was made up of two floating points (an oscillator and a base) that converted the energy from the ocean wave into electrical energy. The RM3 construction dimension detail is indicated in Fig.~\ref{rm3c}, the surface float (oscillator) is oscillating vertically to convert the wave into a mechanical movement. The way to install this device in the ocean is indicated in Fig.~\ref{rm3dep}. the wave surface is causing the oscillators to move against the base, which is fixed by the mooring lines, this device is implemented in the matrix configuration that is many RM3 devices are arranged in column and array to extract more energy from the ocean wave this configuration is indicated in Fig.~\ref{rm3arr}.

\begin{figure}[htbp]
	\centerline{\includegraphics[width=2.5 in]{fig/1.png}}
	\caption{The Construction Detail of The RM3\cite{lawson2014development}.}
	\label{rm3c}
\end{figure}
\begin{figure}[htbp]
	\centerline{\includegraphics[width=2.9 in]{fig/2.png}}
	\caption{The Deployment Configuration of The RM3\cite{neary2014methodology}.}
	\label{rm3dep}
\end{figure}

\begin{figure}[htbp]
	\centerline{\includegraphics[width=2 in]{fig/3.png}}
	\caption{RM3 Array Layout\cite{neary2014methodology}.}
	\label{rm3arr}
\end{figure}
\section{Design of the RM3 with Permanent Magnet linear generator}
According to the manufacturing company \cite{neary2014methodology}, the number of 100 units can generate up to $30 mW$ at $30 kV$ throw rise cable. The design stages for the RM3 device and matrix are the first stage, the location weather such as the privileged wind speed and direction secondly water depths because this device is mounted in the water depth range from $40:100 m$ thirdly the effect of the RM3 matrix in the marine life. Then, determine the operating wave properties. The wave energy flux $Q$ for an irregular type of wave in deep ocean water is shown in Eq.~\ref{J} via "linear wave theory"\cite{reviewofwaveenergy}.

\begin{equation}
Q=\dfrac{\rho \times g^2}{64\pi} \times L^2 \times t\label{J}
\end{equation}


Where the $L$ is the wave height, $t$ is the wave period, $g$ is the gravity acceleration constant, is the $\rho$ water density. More energy can be extracted by the wave that has more height and period. But the device is designed around the point absorber where the maximum energy is extracted or near to the natural frequency of the dominant wave at the deployment site. The power within the sea can be estimated by the Eq.~\ref{P} in terms of the wave high and period.


\begin{equation}
	P(kW/M)=\ 0.49\ L^2 \times t\label{P}
\end{equation}
As previously stated, RM3 is utilised to transfer mechanical energy to electrical energy via one of two techniques, the first is using a hydraulic power take-off system (PTO). The second technique, on the other hand, is less effective in generating power but requires less maintenance because mechanical action is immediately converted into electrical energy using a permanent magnet linear generator (PMLG), Fig.~\ref{rm3lin} depicts the method's design. The linear generator parameters are indicated in Table.~\ref{tab1}.


\begin{figure}[htbp]
	\centerline{\includegraphics[width=1.6 in]{fig/4.png}}
	\caption{The Construction of The RM3 with PM Linear Generator \cite{lawson2014development}.}
	\label{rm3lin}
\end{figure}
\begin{table}[htbp]
	\caption{PM Linear Generator Parameters}
	\begin{center}
			\begin{tabular}{|c|c|c|}
				\hline
				\hline
				\\[-1em]
				Winding resistance&	$4.58$	&$\Omega$ \\
				\hline
				\\[-1em]
				Friction coefficient&	$100$&	--  \\
				\hline
				\\[-1em]
				Magnet pole pitch&	$0.072$&	$m$\\
				\hline
				\\[-1em]
				Flux linkage&	$8$	&$Wb-turn$\\
				\hline
				\\[-1em]
				Winding inductance&	$0.285$	&$H$\\
				\hline
				\hline
			\end{tabular}
		\label{tab1}
	\end{center}
\end{table}
\section{Case study}
Within the scope of Egypt's Vision $2030$, the Egyptian government aimed to make Egypt one of the world's most significant energy hubs by using the country's natural gas reserves in the Mediterranean Sea's east, as well as encouraging the development of renewable energy resources. With $3000\ km$ of coastline, Egypt has a great location. According to the wind speed atlas shown in Fig.~\ref{off} \cite{mortensen2006wind} and the wind speed numerical analysis performed in this study [3], The western Egyptian coastlines, such as Marsa Matruh, Sidi Barrani, and the Gulf of Suez, have a lot of potential for deploying the WEC. With this parameter that is wave hight is $2.5\ m$ and wave period $8\ s$, the extracted power is $24.5\ kW/M$ according to Eq.~\ref{P}.


\begin{figure}[htbp]
	\centerline{\includegraphics[width=3.3 in]{fig/5.png}}
	\caption{Offshore Wind Speed Atlas of Egypt \cite{mortensen2006wind}.}
	\label{off}
\end{figure}

\subsection{Result of the wave simulation}
The ocean wave simulation is performed using the WecSim toolbox \cite{lucas2012development} with irregular wave type (the real wave action) and $L=2.5 m$ and $t=8 s$ the resultant simulation is indicated in this section, firstly the Pierson-Moskowitz Spectrum that is the distribution of the wave energy (the frequency and direction) is indicated in Fig.~\ref{pm}. The wave surface shape is indicated in Fig.~\ref{WSE} over 1000 seconds. The two bodies responses for the wave surface Elevation are indicated in Fig.~\ref{FaT}, the forces in this figure are completely articulated in this chapter \cite{pecher2017erratum}, the used force to generate useful work (electrical energy in this case) is the excitation force which is marked in orange color, from this figure, the excitation force on the oscillator is the inverse in phase with that acting in the base that is because the oscillator is transferring the force with small resistant from the base. In this case, if the RM3 is equipped with PM linear generator then this will extract this excitation force.


\begin{figure}[htbp]
	\centerline{\includegraphics[width=2.9 in]{fig/6.png}}
	\caption{Pierson-Moskowit Spectrum at $L=2.5\ m$ and $t=8\ s$}
	\label{pm}
\end{figure}
\begin{figure}[htbp]
	\centerline{\includegraphics[width=3.3 in]{fig/7.png}}
	\caption{Wave Surface Elevation}
	\label{WSE}
\end{figure}
\begin{figure}[htbp]
		\centerline{\includegraphics[width=3.3 in]{fig/8.png}}
	\caption{ Forces Acts on The Two Bodies of The RM3 }
	\label{FaT}
\end{figure}
\subsection{Result of the RM3 with PM linear generator}
In this section, the output power from the linear generator is demonstrated in two figures the upper figure contains the full time simulation while the lower figure is a sectional zoom of the full simulation time. the PTO force and friction force are indicated in Fig.~\ref{Frfr}, the PTO force is the used force to generate the Electrical energy, Fig.~\ref{Abpwr} shows the absorbed power and electrical power generated from the linear generator. The output voltage of the pm linear generator is shown in Fig.~\ref{vpm}. The next stage is using the power electrical converters and storage systems to regulate the output voltage to use in the home and industrial applications.


\begin{figure}[htbp]
	\centerline{\includegraphics[width=3.3 in]{fig/9.png}}
	\caption{ Friction Force and PTO Force  }
	\label{Frfr}
\end{figure}
\begin{figure}[htbp]
	\centerline{\includegraphics[width=3.3 in]{fig/10.png}}
	\caption{ Electrical and Absorbed Power }
	\label{Abpwr}
\end{figure}
\begin{figure}[htbp]
	\centerline{\includegraphics[width=3.3 in]{fig/11.png}}
	\caption{ The PMLG's Output Voltage }
	\label{vpm}
\end{figure}

\begin{figure*}[ht]
	\centerline{\includegraphics[width=4.3 in]{fig/12.png}}
	\caption{ System Description  }
	\label{sys}
\end{figure*}
\begin{figure}[htbp]
	\centerline{\includegraphics[width=2.9 in]{fig/13.png}}
	\caption{ Electrical Equivalent Circuit Model}
	\label{Electrical Equivalent Circuit Model}
\end{figure}
\begin{figure}[htbp]
	\centerline{\includegraphics[width=2 in]{fig/14.png}}
	\caption{ Basic Configuration of DC-DC Converter}
	\label{Basic Configuration of DC-DC Converter}
\end{figure}

\section{Description of Wave-Energy Power Generation With Battery-Based Charging System}
\subsection{System configuration}
In the proposed system, the PM linear generator acts a primary power source and battery is used as a secondary storage unit as shown in Fig.~\ref{sys}. Due to the irregular moving of the waves, the most factors which should be taken into considerations are the wave displacement of two crests (which affecting the voltage frequency) and the speed/height of the wave (which affecting the voltage amplitude). Moreover, the frequency and magnitude of the linear-generator voltage have a wide variation range.


Energy storage systems (ESSs) appear to be effective solutions for tackling the high fluctuation of electricity generating variability and power quality supplied to the grid. This returns to the merits of ESSs to give a high level of energy in a short period and also can supply the electrical network with energy for a long time.


The ability to react very quickly and cope with many cycles of discharge/charge per day are the most important requirements for ESSs applied for both applications of fluctuations mitigation in voltage and installations of waves.


There are different classifications of ESSs depending on the various shares of markets and specifications. However, one of the ways to increase the output power of wave energy systems is using high power buck-boost DC converter, that converter is used to charge/discharge the battery storage unit according to the presented management rules.
%
\subsection{Modeling of battery}
In this study, the whole performance of battery with all characteristics of current, voltage, SOC, and temperature is attracting a great attention rather than the electromechanical behavior.  In this work, the stated equivalent circuit model of a battery is used. The detailed electrical equivalent circuit model of the used battery is shown in Fig.~\ref{Electrical Equivalent Circuit Model}.
\subsubsection{State of charge calculation}
\begin{equation}
	soc=soc_0 \times e^{-5.738(A-B)}-\frac{1}{3600\times1.1}\times \int_{t1}^{t2}{i_b\ dt}\label{Soc}
\end{equation}

Where, $A$ is $\frac{1}{T-209.9}$, and $B$ is $\frac{1}{T_{ref}-209.9}$, $SOC_0$ is the initial $SOC$, $T$ is the cell Temperature, $T_{ref}$ is the reference Temperature, $i_{b}$ is the battery cell current, and $t$ is the time.
\subsubsection{Calculation of voltage drops}
\begin{equation}
	V_0 = i_b \times R_0\label{vo}
\end{equation}
\begin{equation}
	V_1=\int_{0}^{t}{\left(\frac{i_b}{C_1}\right)-\left(\frac{V_1}{R_1\times C_{1}}\right) dt} \label{v1}
\end{equation}
\begin{equation}
	V_2=\int_{0}^{t}{\left(\frac{i_b}{C_2}\right)-\left(\frac{V_2}{R_2\times C_{2}}\right) dt} \label{v2}
\end{equation}
\begin{equation}
	V_{cell}= V_{oc}-V_0-V_1-V_2\label{vcell}
\end{equation}
\begin{equation}
	i_b = I_b/np\label{ib}
\end{equation}
\subsection{Modeling of DC-DC converter}
As shown in Fig.~\ref{sys}, the DC-DC converters, illustrated in Fig.~\ref{Basic Configuration of DC-DC Converter}, are required in the presented system whose design details can be given in details as follows.


\textbf{The common components of the DC-DC circuit:}

\begin{enumerate}
	\item A detection part is required to predict the voltage level and take the appropriate action for bucking or boosting.
	\item The switches are given as the MOSFETs which are preferred for the converters of high frequency.
	\item Diodes shown as Shockley types, These diodes can function at high frequencies since they have a low forward junction voltage while conducting.
	\item Inductor with appropriate value to keep current within certain range and prevent current from reaching zero value.
	\item Capacitor with appropriate value to keep voltage within certain range and prevent voltage from reaching zero value.
	\item A sensor that senses the input and output value to make a controller.
\end{enumerate}
The output voltage is given as:
\begin{equation}
	v_o=\ \frac{v_s}{(1-D)}\label{V0}
\end{equation}
\begin{equation}
	L_{min}=\frac{\left(1-D\right)^2 \times R \times D}{2f},\ L>L_{min}\label{lmin}
\end{equation}
\begin{equation}
	{C_{min}=\frac{D\times v_o}{\Delta v\times R \times f}},\ C>C_{min}\label{cmin}
\end{equation}

where, $v_o$ is the output voltage, $v_s$ is the input voltage, $D$ is the duty cycle, $R$  represents the load,(dc bus), $f$ is the switching frequency. $\Delta v$ is the value of the ripple voltage, $L_{min}$ is the minimum value of the inductance,and $C_{min}$ is the minimum value of the capacitance.

\section{PERFORMANCE TEST OF THE OVERALL SYSTEM}
In this section, some of the results are given to test the performance of the adopted wave energy system with the energy storage and the terminal load side as shown in Fig. 12. The test is performed at regular wave with $L = 2.5 m$ and $t = 8 s$ with a desired output boost voltage set as 260 V. Fig.~\ref{Performance of overall system} shows the response of the overall control system including the three-phase generated voltage from the PM linear generator, the rectified output voltage, the output from boost converter, and the electrical power in the load side. In addition, the output DC current of load is given in Fig.~\ref{DC ccurrent at load}.


\begin{figure}[htbp]
	
	\subfloat{%
		\includegraphics[clip,width=\columnwidth]{fig/15-a.png}%
	}

	\subfloat{%
		\includegraphics[clip,width=\columnwidth]{fig/15-b.png}%
	}

	\subfloat{%
	\includegraphics[clip,width=\columnwidth]{fig/15-c.png}%
	}

	\subfloat{%
	\includegraphics[clip,width=\columnwidth]{fig/15-d.png}%
	}

	\caption{Performance of Overall System}
	\label{Performance of overall system}
	
\end{figure}

\begin{figure}[htbp]
	
	\subfloat{%
		\includegraphics[clip,width=\columnwidth]{fig/16-a.png}%
	}
	
	\subfloat{%
		\includegraphics[clip,width=\columnwidth]{fig/16-b.png}%
	}
	
	
	\caption{ DC Current at Load}
	\label{DC ccurrent at load}
	
\end{figure}

\begin{figure}[htbp]
	
	\subfloat{%
		\includegraphics[clip,width=\columnwidth]{fig/17-a.png}%
	}
	
	\subfloat{%
		\includegraphics[clip,width=\columnwidth]{fig/17-b.png}%
	}
	
	
	\caption{The Effect on The Battery Performance without Using Super Capacitor.}
	\label{The effect on the battery performance without using super capacitor.}
	
\end{figure}

\begin{figure}[htbp]
	
	\subfloat{%
		\includegraphics[clip,width=\columnwidth]{fig/18-a.png}%
	}
	
	\subfloat{%
		\includegraphics[clip,width=\columnwidth]{fig/18-b.png}%
	}
	
	
	\caption{Battery Performance when Using Supercapacitor}
	\label{Battery performance when using supercapacitor}
	
\end{figure}

 
Moreover, the battery performance is shown in Fig.~\ref{The effect on the battery performance without using super capacitor.} including the current and voltage. The transient of battery performance can be improved using a supercapacitor connected with the load side \cite{last}. The performance of the battery after using the supercapacitor is illustrated in Fig.~\ref{Battery performance when using supercapacitor} which assures the good behavior of the adopted system.

\section{Conclusions}

Renewable energy sources are the latest global concern in the generation of electricity due to the numerous problems of fossil fuels. This paper has successfully built a model that simulate a wave Energy Converter to continuously supply a DC load at a good power quality. In addition, the Buck-Boost Bi-direction converter has been utilized to decrease the wave power availability problem. The results have assured the idea that the small remote areas can depend on the wave energy generated power with optimized control and management system for the hybrid energy storage (battery and super-capacitor).

\bibliographystyle{./IEEEtran}
\bibliography{./IEEEabrv,./IEEEexample}
\end{document}
